\section{System Model and Methodology}

\subsection{System Overview}

The proposed system consists of a hybrid energy architecture integrating solar photovoltaic (PV) generation, battery energy storage, grid supply, diesel backup generation, and electric vehicle (EV) charging demand. The system is designed for weak-grid environments such as those observed in parts of Kenya, where grid reliability is not guaranteed.

The objective of the system is to ensure continuous EV charging while minimizing operational cost, diesel fuel consumption, and carbon emissions.

---

\subsection{Solar Photovoltaic Power Model}

The output power from the solar PV system is modeled as:

\begin{equation}
P_{pv}(t) = \eta_{pv} \cdot A \cdot G(t)
\end{equation}

where:
\begin{itemize}
    \item $\eta_{pv}$ is the PV conversion efficiency
    \item $A$ is the panel area
    \item $G(t)$ is solar irradiance at time $t$
\end{itemize}

This model captures the time-varying nature of solar energy availability.

---

\subsection{Battery Energy Storage Model}

The state of charge (SOC) of the battery is governed by:

\begin{equation}
SOC(t+1) = SOC(t) + \frac{P_{ch}(t)\Delta t}{E_{cap}} - \frac{P_{dis}(t)\Delta t}{E_{cap}}
\end{equation}

subject to:

\begin{equation}
0 \leq SOC(t) \leq 1
\end{equation}

where:
\begin{itemize}
    \item $E_{cap}$ is battery energy capacity
    \item $P_{ch}(t)$ is charging power
    \item $P_{dis}(t)$ is discharging power
\end{itemize}

---

\subsection{EV Charging Demand Model}

EV charging demand is modeled as a time-varying load:

\begin{equation}
P_{ev}(t) = \sum_{i=1}^{N} P_i(t)
\end{equation}

where $N$ represents the number of EVs connected to the charging station.

---

\subsection{Grid Power Model under Weak Conditions}

Grid power availability is constrained as:

\begin{equation}
P_{grid}(t) \leq P_{grid}^{max}(t)
\end{equation}

where $P_{grid}^{max}(t)$ varies due to outages and instability, representing weak-grid behavior observed in Kenyan distribution networks.

---

\subsection{Diesel Generator Fuel Consumption Model}

Diesel generator fuel consumption is modeled as:

\begin{equation}
F(t) = \alpha \cdot P_{diesel}(t)
\end{equation}

where:
\begin{itemize}
    \item $\alpha$ is the fuel consumption coefficient
    \item $P_{diesel}(t)$ is diesel generator output power
\end{itemize}

---

\subsection{Carbon Emission Model}

Total emissions are calculated as:

\begin{equation}
E_{CO2}(t) = \beta \cdot P_{diesel}(t)
\end{equation}

where $\beta$ is the emission factor per unit diesel-generated energy.

---

\subsection{Energy Balance Constraint}

At every time step, the system must satisfy:

\begin{equation}
P_{pv}(t) + P_{grid}(t) + P_{dis}(t) + P_{diesel}(t) = P_{ev}(t) + P_{ch}(t)
\end{equation}

---

\subsection{Energy Management Objective}

The system minimizes a weighted cost function:

\begin{equation}
\min \; J = w_1 C_{fuel} + w_2 E_{CO2} + w_3 P_{grid}^{stress}
\end{equation}

where:
\begin{itemize}
    \item $C_{fuel}$ represents diesel fuel cost
    \item $E_{CO2}$ represents emissions
    \item $P_{grid}^{stress}$ penalizes grid overloading
    \item $w_1, w_2, w_3$ are weighting factors
\end{itemize}